\RequirePackage{iftex}
\ifPDFTeX
\documentclass[a4paper,11pt]{article}
\usepackage[hmargin=14mm,vmargin=14mm]{geometry}
\usepackage[T1]{fontenc}
\usepackage[utf8]{inputenc}
\usepackage{CJKutf8}
\else
\documentclass[a4paper,11pt]{jsarticle}
\usepackage[dvipdfmx,hmargin=14mm,vmargin=14mm]{geometry}
\fi
\usepackage{amsmath,amssymb}
\usepackage{multicol}

\setlength{\parindent}{0pt}
\setlength{\parskip}{0.35em}
\setlength{\columnsep}{7mm}
\setlength{\multicolsep}{0.5em}

\begin{document}
\ifPDFTeX
\begin{CJK}{UTF8}{min}
\fi

\newcommand{\problem}[2]{\par\noindent\mbox{問題 #1（#2点）}\enspace\ignorespaces}

\begin{center}
{\Large 解析学 II 中間試験}\\[0.4em]
\end{center}

全25問、各10点、合計250点。途中の計算や判断の理由も書きなさい。

\begin{multicols}{2}
\raggedcolumns
\problem{1}{10}
$x>0$ に対して、$\arctan x+\arctan(1/x)$ が $x$ によらず一定であることを示し、その値を求めよ。

\problem{2}{10}
分母を平方完成し、次の不定積分を求めよ。
$$
\int\frac{x+3}{x^2+2x+5}\,dx
$$

\problem{3}{10}
長方形 $K=[0,2]\times[0,1]$ の各辺を $n$ 等分する。各小長方形の右上の頂点を代表点として、$f(x,y)=xy$ のリーマン和を $n$ の式で表せ。その極限から $\iint_K xy\,dx\,dy$ を求めよ。

\problem{4}{10}
次の積分を求めよ。
$$
\int_0^1\left(\int_x^1 e^{y^2}\,dy\right)dx
$$

\problem{5}{10}
広義積分 $\displaystyle\int_1^\infty\frac{dx}{(x+1)^{3/2}}$ の値を求めよ。

\problem{6}{10}
$D=\{(x,y):0\le x+y\le 2,\ 0\le x-y\le 3\}$ とする。$u=x+y$, $v=x-y$ とおき、$x,y$ を $u,v$ の式で表せ。この変換を $T(u,v)=(x,y)$ とおき、ヤコビアン $J_T=\frac{\partial(x,y)}{\partial(u,v)}$ を求め、それを用いて $\iint_D(x+y)(x-y)\,dx\,dy$ を計算せよ。

\problem{7}{10}
立体 $V=\{(x,y,z):x^2+y^2\le 1,\ 0\le z\le 2-x^2-y^2\}$ の体積を求めよ。

\problem{8}{10}
上半球 $H=\{(x,y,z):x^2+y^2+z^2\le 4,\ z\ge 0\}$ の密度を $\rho(x,y,z)=x^2+y^2+z^2$ とする。球面座標を用いて、積分範囲と体積要素を示し、質量を求めよ。

\problem{9}{10}
$xy$ 平面上の領域 $D=\{(x,y):x,y>0,\ 1\le xy\le 4,\ 1\le x/y\le 4\}$ を考える。
\par (1)\enspace 領域 $D$ を $xy$ 平面上に図示せよ。
\par (2)\enspace $I=\iint_D(x^2+y^2)\,dx\,dy$ を求めよ。

\problem{10}{10}
小問ごとの配点は可変とする。\par
半径 $R>0$ の円形のピザにチーズが偏って載っている。ピザの領域を $D=\{(x,y):x^2+y^2\le R^2\}$ とし、チーズの面密度を $D$ 上の連続な非負関数 $f(x,y)$ とする。チーズの総質量を $M>0$ とする。ピザの中心を通り、偏角 $\theta$ の半直線を含む直線で切ったとき、偏角が $\theta$ から $\theta+\pi$ までの半円に載ったチーズの質量を $H(\theta)$ とする。$H(\theta)$ が連続であることは証明なしで用いてよい。
\par (1)\enspace 極座標 $x=r\cos\varphi$、$y=r\sin\varphi$ を用いて、$M$ と $H(\theta)$ を定積分で表せ。
\par (2)\enspace ピザの面積とチーズの質量を同時に二等分する直線が必ず存在することを示せ。
\par (3)\enspace $f(x,y)=\rho_0\left(1+\frac{x}{2R}+\frac{y}{3R}\right)$、$\rho_0>0$ のとき、二等分する直線の方程式を求めよ。

\columnbreak

\problem{11}{10}
$f(x)=(2x-1)/(x+1)$（$x>-1$）の逆関数を、その定義域とともに求めよ。

\problem{12}{10}
部分分数分解を用いて次の積分を求めよ。
$$
\int_0^1\frac{dx}{(x+1)(x+2)}
$$

\problem{13}{10}
$D=[-1,1]\times[0,2]$ とする。対称性を用いて $\iint_D(x^3y+y^2)\,dx\,dy$ を求めよ。

\problem{14}{10}
$K=[0,2]\times[0,1]$ の各辺を二等分する。各小長方形の右上の頂点を代表点としたときの $f(x,y)=x+y$ のリーマン和を求めよ。

\problem{15}{10}
曲線 $y=x^2$ と直線 $y=2x$ に囲まれた領域の面積を重積分で求めよ。

\problem{16}{10}
次の積分を、積分順序を交換した式で表せ。
$$
\int_0^1\left(\int_{\sqrt{y}}^1 x^2\,dx\right)dy
$$

\problem{17}{10}
広義積分 $\displaystyle\int_0^1\frac{dx}{\sqrt{1-x}}$ の値を求めよ。

\problem{18}{10}
$Q=[1,\infty)^2$ とする。$[1,L]^2$ 上の積分の $L\to\infty$ での極限として、広義積分 $\iint_Q x^{-2}y^{-2}\,dx\,dy$ を求めよ。

\problem{19}{10}
$D=\{(x,y):1\le x\le2,\ x\le y\le3x\}$ とする。$u=x$、$v=y/x$ とおいて $\iint_D(y/x)\,dx\,dy$ を求めよ。

\problem{20}{10}
$D=\{(x,y):1\le x^2+y^2\le4,\ x\ge0,\ y\ge0\}$ とする。極座標で $\iint_D(x^2+y^2)\,dx\,dy$ を求めよ。

\problem{21}{10}
変数変換 $x=u$、$y=2v$ を用いて、楕円板 $4x^2+y^2\le4$ の面積を求めよ。

\problem{22}{10}
$x,y,z\ge0$、$x+y+z\le2$ で定まる立体の体積を重積分で求めよ。

\problem{23}{10}
円柱座標を用いて、$x^2+y^2\le4$、$0\le z\le2-\sqrt{x^2+y^2}$ で定まる円すいの体積を求めよ。

\problem{24}{10}
球面座標を用いて、$1\le x^2+y^2+z^2\le4$ で定まる立体の体積を求めよ。

\problem{25}{10}
円柱 $x^2+y^2\le1$、$0\le z\le2$ の密度が $1+x^2+y^2$ のとき、この円柱の質量を求めよ。
\end{multicols}

\ifPDFTeX
\end{CJK}
\fi

\end{document}
